\section{系统视角：吞吐受限时先查哪里}
对于 Agentic RL 训练，推理模式与显存占用并不是两个独立话题。真正的系统瓶颈往往体现在：一边想提高 rollout 吞吐，一边又被 context、batch 和 cache 挤爆显存。理解 online/offline 的差异，本质上是在理解训练系统如何在延迟和吞吐之间做取舍。

\begin{knowledgebox}{排查顺序建议}
\begin{itemize}
  \item 先看每轮 rollout 的 token 数和并发数
  \item 再看 KV cache、batch size 是否占满显存
  \item 最后再决定采用 online 还是 offline 推理编排
\end{itemize}
\end{knowledgebox}

\subsection{本章小结}
显存分析不是附属话题，而是 RL 系统设计的一部分。只有吞吐、延迟和显存三者同时平衡，训练管线才能稳定跑起来。

\newpage
